\usemodule[visual]

\setscript[hanzi]
\setupalign[hanging, hz]

\usemodule[zhfontsext]
\usetypescript[cnfontb]
\setupbodyfont[cnfontb,rm,12pt] % 小四号 = 12pt
\mainlanguage[cn]

\setupinterlinespace[line=3.66ex] %line=3.05ex相当于WORD 中的1.0倍行距, 3.66ex相当于WORD 中的1.2倍行距

\setuppapersize[A4]
\setuplayout[width=fit,height=fit,backspace=10mm,cutspace=10mm,topspace=10mm,margin=10mm,header=10mm,footer=10mm]
\setuppagenumbering[location={header,right},style=\bfb]

\define[1]\exampleLine{
    \blank[10pt]
    \framed[frame=off,align=flushleft,offset=3pt,strut=none,width=\textwidth,height=fit,background=color,backgroundcolor=gray]{{\bfx #1}}
    \blank[5pt]
}


\startbuffer[fruita]
\starttabulate[format=|c|c|c|]
\HL
\NC 年份 \NC 产品 \NC 销量 \NC\NR
\HL
\NC 2021 \NC 柚子 \NC 4500 \NC\NR
\NC 2022 \NC 水蜜桃 \NC 480 \NC\NR
\HL
\stoptabulate
\stopbuffer

\starttext


\startbuffer[fruitb]
\starttabulate[format=|c|c|c|]
\HL
\NC 年份 \NC 产品 \NC 销量 \NC\NR
\HL
\NC 2023 \NC 柠檬 \NC 50 \NC\NR
\NC 2024 \NC 蓝莓 \NC 930 \NC\NR
\HL
\stoptabulate
\stopbuffer

% % % ================== 代码示例 - 1
\fakewords{60}{80}

\startplacetable [location={none}]     % 设置最外层不显示表题
  \startfloatcombination [nx=2,ny=1]
\startplacetable[title=fruit A]
{
\getbuffer[fruita]
}
\stopplacetable
\startplacetable[title=fruit B]
{
\getbuffer[fruitb]
}
\stopplacetable
\stopfloatcombination
\stopplacetable

\fakewords{60}{80}

\page

% % % ================== 代码示例 - 2
\fakewords{60}{80}

\startplacetable[title=共享表序表题]
  \startcombination [nx=2,ny=1]
\startcontent
\getbuffer[fruita]
\stopcontent
\startcaption
fruit A
\stopcaption

\startcontent
\getbuffer[fruitb]
\stopcontent
\startcaption
fruit B
\stopcaption

\stopcombination
\stopplacetable

\fakewords{60}{80}

\stoptext